% Draft subsection for Part III, to follow \subsection{Diagnosing a negative result}
% in draft6.tex. Status 2026-08-02: numbers are from PILOT/discovery data
% (environment seeds 901--903); the registered confirmation has NOT been run and the
% registration is held at :design-review-hold. Nothing here should be stated as a
% confirmed result until that changes.

\subsection{Running the diagnostic}
\label{sec:diagnostic-run}

The diagnostic of the previous subsection was run against an evolved genome that
produced no certificate. Its purpose is narrow, and worth restating plainly before
the numbers: it does not search, it does not exhibit a path, and it cannot
establish assimilation. It asks one question at a time, at a single locus, about a
genome that already exists.

The design is a panel of sixteen loci, stratified by position and by endpoint
density, each evaluated under six conditions. Two conditions leave the inherited
rule in place, free to rewrite or held. Two install a rule that an earlier
endpoint-mapping run recommended for that locus, again free or held. One installs a
rule matched to the recommended one in Hamming distance but not recommended, held.
The last takes the recommended rule, holds it, and evaluates it on rewrite tapes
that played no part in building the map. Each condition is scored on the band score
of \secref{sec:endogenous-gain}, paired across environment seed, rewrite tape and
evaluation site.

The three cases of the previous subsection map onto contrasts among these
conditions. Whether \emph{any} static rule preserves function at a locus is read
from the held conditions against threshold. Whether the difficulty is that mutation
cannot prepare the allele is read from the recommended rule against the inherited
one, both held. Whether an assimilation step was simply available and not taken is
read from the inherited rule held against the same rule left plastic.

On the pilot panel the third case fires. Holding the inherited allele does not cost
function: mean band rises from $0.181$ plastic to $0.196$ held. The inherited rule
therefore already carries its behaviour without within-evaluation rewriting at these
loci, which by the previous subsection is a statement about the search rather than
about the substrate --- an assimilation step was accessible and was not taken.

The second case fires as well, and more strongly. The recommended rule beats the
inherited one by $0.031$ in mean band when both are held ($0.227$ against $0.196$),
and beats a Hamming-matched unrecommended rule by $0.031$ ($0.227$ against $0.196$).
The same ordering appears without holding, at $0.219$ against $0.181$. So the
endpoint map is carrying real information about which rule belongs at a locus, and
that information is visible whether or not the locus is frozen. Held on rewrite
tapes it never saw, the recommended rule retains $0.223$, so the map is not an
artefact of the tapes it was built from.

Two remarks on what these numbers do \emph{not} say.

First, the effect of holding is an order of magnitude smaller than the effect of
content ($0.015$ against $0.031$), and it does not discriminate between rules: a
held unrecommended rule performs as well as a held inherited one ($0.196$ in both
cases). The apparatus is detecting that some rules are better than others, not that
freezing is what makes them work.

Second, and against the natural reading of the certificate, realized rewriting is
\emph{not} a usable secondary signal here. Counting rule changes actually made, a
held arm does less rewriting than its plastic counterpart by $132$ and $121$ changes
respectively, against a ceiling of $120$ --- the number of steps in an evaluation,
and hence the most that freezing a single locus can remove. The reduction is
therefore almost exactly the held locus's own rewrites and nothing further; a held
arm satisfies ``did less work'' by construction, at any rule, including one chosen
to be bad. A criterion combining preserved function with reduced work would be met
tautologically at this intervention scale, and the same objection applies to any
cost term that is a capacity rather than a record of work performed.

That second remark bears directly on \eqref{eq:selectable}. Writing fitness as
$\mathrm{perf} - c\,\mathrm{dep}$ makes the cost term flat in behaviour whenever
$\mathrm{dep}$ is a fraction over unheld loci: holding one locus of $W$ then shifts
fitness by a constant $c/W$ regardless of what the held rule does. At the scale of a
single-locus intervention that constant is small in absolute terms and decisive in
relative ones, because it applies identically to every paired comparison and so
resolves every near-tie in the same direction. Contrasts computed on
$\mathrm{fit}_c$ can consequently report a uniform advantage for holding that
survives no restatement on $\mathrm{perf}$ alone. The remedy is not a different $c$;
it is to decide the comparison on the behavioural score and to report work
separately, which is what is done above. We record this because the failure is
invisible in the trajectory, is not detected by any invariant of
\secref{sec:apparatus}, and was found only by recomputing the same contrasts on the
score's components.

The panel is a diagnostic and is treated as discovery data. Its sixteen loci are
also its inferential units: a within-locus sign rule collapses more than a thousand
paired evaluations per condition into sixteen decisions, and a mean difference of
$0.031$ that is plainly present in the pooled data need not clear a familywise
threshold stated over those sixteen. Whether the effects reported here survive as
registered results is a question for a confirmation run on held-out environment
seeds, under a criterion stated on the behavioural score, and that run has not been
performed.
